\section{Introduction}
\label{sec:introduction}

The widespread adoption of cloud-native computing and container orchestration platforms, led by Kubernetes, has transformed contemporary software systems into complex distributed webs of loosely coupled microservices \cite{gannon_2017, bogner2020microservices}. While this architectural paradigm accelerates feature delivery, independent scalability, and organizational agility, it fundamentally complicates operational visibility. When transient network anomalies, CPU saturation, or pod failures occur within polyglot service graphs, failure symptoms cascade across remote procedure calls (RPC) and event brokers, obfuscating the actual root cause \cite{soldani2022anomaly, zhou2021fault}.

To diagnose such failures, modern DevOps and Site Reliability Engineering (SRE) teams collect the ``Three Pillars of Observability'': numerical time-series metrics, structured event logs, and distributed trace spans \cite{sridharan_2018, beyer2016sre}. However, in current industrial practice and academic literature, these telemetry modalities are frequently managed in fragmented silos \cite{kosinska_2023, nedelkoski2020selfattentive}. SREs are forced to manually correlate dashboard anomalies in Prometheus with unstructured text logs in Grafana Loki and call trees in Grafana Tempo or Jaeger. This manual triaging prolongs the Mean Time to Diagnose ($MTTDiag$), heightens cognitive load, and frequently leads to erroneous incident remediation.

Furthermore, naive automated Root Cause Analysis (RCA) approaches that rely on a single telemetry modality suffer from severe architectural limitations:
\begin{itemize}[leftmargin=1.2em, itemsep=2pt, topsep=3pt]
    \item \textbf{Cascading Span Bias in Monomodal Tracing:} Naive trace diagnostic algorithms that isolate culprits by ranking inclusive span durations invariably point to upstream ingress microservices (e.g., API gateways or edge proxies), misidentifying victims of accumulated downstream latency as the source of failure \cite{yu2021microrank}.
    \item \textbf{Semantic Blindness in Monomodal Metrics:} System-level metric scrapers capture host resource spikes (e.g., CPU, RAM) but lack transaction context, making them incapable of distinguishing between legitimate user traffic surges and rogue background processes.
    \item \textbf{Invasiveness and Observability Tax:} Increasing probe frequency and capturing un-sampled traces across high-throughput clusters introduces non-negligible CPU and memory overhead, perturbing application performance and risking the very stability observability seeks to ensure \cite{vancamp24, valbusa_2025}.
\end{itemize}

Despite these pressing dilemmas, the software engineering community lacks a formal, empirically validated reference architecture that defines how multimodal telemetry streams should be semantically unified, correlated, and bounded in resource footprint. 

To bridge this gap, this paper introduces \textbf{PANOPTES}, a novel reference architecture and multimodal correlation engine for cloud-native microservice observability. Grounded in the epistemology of \textbf{Design Science Research (DSR)} \cite{simon_1969, hevner_2004, peffers_2007}, PANOPTES delivers a 4-layer reference model that unifies non-invasive eBPF and OpenTelemetry probes, federated multimodal storage backends, and a causal correlation engine powered by intra-trace \textit{Span Self-Time Decomposition} ($T_{\text{self}}$) and OpenMetrics Exemplar binding.

Specifically, this research addresses three core Research Questions (RQs):
\begin{itemize}[leftmargin=1.2em, itemsep=2pt, topsep=3pt]
    \item \textbf{RQ1 (Architectural Design \& Semantic Invariants):} What structural layers, semantic bindings, and design heuristics are necessary to eliminate telemetry silos and prevent Cascading Span Bias in distributed microservices?
    \item \textbf{RQ2 (Diagnostic Efficacy \& Latency Reduction):} How effectively does the PANOPTES multimodal correlation engine isolate root causes under cascading chaos perturbations compared to native Kubernetes inspection and monomodal distributed tracing?
    \item \textbf{RQ3 (Resource Overhead \& Invasiveness Envelope):} What is the steady-state CPU and memory footprint imposed by the complete PANOPTES reference stack, and can it be bounded within a predictable operational envelope?
\end{itemize}

Through a rigorous controlled experiment over the polyglot CNCF OpenTelemetry Astronomy Shop benchmark ($N=30$ paired trials, stratified across network delay, pod termination, and CPU stress), we demonstrate that PANOPTES achieves $\mathbf{83.33\%}$ diagnostic accuracy, outperforming monomodal tracing ($0.00\%$, $p < 10^{-11}$) and native CLI inspection ($20.00\%$, $p < 10^{-6}$), while reducing median $MTTDiag$ from $300.00\,\text{s}$ to $\mathbf{17.27\,\text{s}}$ and bounding observability CPU overhead to under $10\%$ of cluster capacity.
